% Example response (revision) letter: a cover letter to the editors, then one
% section per reviewer. Each reviewer's comments are quoted verbatim in that
% reviewer's color, numbered by page or section as in the report; each reply
% follows in black. All names, addresses and comments below are placeholders.
\documentclass[sigleft]{letter}
\usepackage[hidelinks]{hyperref}
\usepackage{microtype}
\usepackage{amsmath}
\usepackage{bm}
\usepackage{xcolor}
\usepackage[a4paper,left=2.5cm,right=2.5cm,top=2.5cm,bottom=2.5cm]{geometry}

% The letter class has no sections: define \section and \subsection.
\newcounter{section}
\newcounter{subsection}[section]
\setcounter{secnumdepth}{3}

\makeatletter
% Opening: place and date on the left, sender's address on the right.
\renewcommand*{\opening}[1]{\ifx\@empty\fromaddress%
  \thispagestyle{firstpage}%
    {\raggedleft\@date\par}%
  \else%
   \thispagestyle{empty}%
   {%
    \begin{minipage}[c]{0.70\linewidth}
    \vspace{-3cm}
    City, \@date\\[2ex]
    \toname \\
    \toaddress
    \end{minipage}
    \begin{minipage}[c]{0.25\linewidth}
    \raggedleft\begin{tabular}{l@{}}\ignorespaces
    \fromaddress\\[0.0em]
    \end{tabular}
    \end{minipage}
    \par
}%
  \fi
  \vspace{2\parskip}%
  #1\par\nobreak}

\newcommand\section{\@startsection{section}{1}{\z@}%
                                  {-3.5ex \@plus -1ex \@minus -.2ex}%
                                  {2.3ex \@plus.2ex}%
                                  {\normalfont\Large\bfseries}}
\newcommand\subsection{\@startsection{subsection}{2}{\z@}%
                                     {-3.25ex\@plus -1ex \@minus -.2ex}%
                                     {1.5ex \@plus .2ex}%
                                     {\normalfont\large\bfseries}}
\renewcommand\thesection{\@arabic\c@section}
\renewcommand\thesubsection{\thesection.\@arabic\c@subsection}
\makeatother

% One color per reviewer, used for the section title and the quoted comments.
% The color of \reviewerA is also the one used in the manuscript for changes.
\newcommand{\reviewerA}[1]{{\color{blue}#1}}
\newcommand{\reviewerB}[1]{{\color{red}#1}}
\newenvironment{quoteA}{\begin{quote}\itshape\color{blue}}{\end{quote}}
\newenvironment{quoteB}{\begin{quote}\itshape\color{red}}{\end{quote}}
% Replies, nested inside the reviewer's quote environment.
\newenvironment{reply}{\begin{quote}\upshape\color{black}}{\end{quote}}
% Placeholders for replies still waiting for the authors.
\newcommand{\TODO}[1]{\textcolor{orange}{\textbf{[TODO: #1]}}}

\signature{Corresponding Author\\on behalf of the co-authors,\\Second Author and\\Third Author.\\[2ex]}
\address{ }

\begin{document}

\begin{letter}{}
\opening{Dear Editor,}

Subject: submission of the revised manuscript ``Title of the manuscript''.

We are submitting a revised version of the manuscript ``\textit{Title of the manuscript}''.
We addressed the reviewers' suggestions and revised the manuscript to improve its clarity and completeness.
In particular, we [one or two sentences summarising the main changes].

We attach a letter of response to the reviewers describing the modifications we made, with the changes to the manuscript highlighted in \reviewerA{blue}.

We thank you and the reviewers for the time and attention devoted to our manuscript.

\closing{Sincerely yours,}

\noindent
\emph{Corresponding author's address}:

\noindent
Corresponding Author\\
Department, Institution\\
Street, Postcode, City, Country\\
E-mail: author@example.org

\end{letter}

\newpage

\reviewerA{\section{Reviewer \#1}}

We thank the reviewer for the time spent reading and commenting on the manuscript. We start by quoting the general comment,

\begin{quoteA}
``The paper proposes a reduced-order model for parametrised problems on unstructured meshes. The method is interesting and the paper is well written. I recommend publication after the minor issues below are addressed.''
\end{quoteA}

We address the specific comments below.
\begin{quoteA}
Specific comments:
\begin{enumerate}
\item[p. 4] The update rule in Equation (2) adds the time derivative to the state without multiplying it by $\Delta t$. Shouldn't it be $\bm u_{k+1} = \bm u_k + \Delta t\, f(\bm u_k)$?
\begin{reply}
We thank the reviewer for spotting this inconsistency: with our notation, the update indeed reads $\bm u_{k+1} = \bm u_k + \Delta t\, f(\bm u_k)$, since the network learns the time derivative itself.
We corrected Equation (2) in Subsection 1.1.
% Keep the new text of the manuscript here, commented, for the co-authors:
% ``... allowing the solution to be updated as $\bm u_{k+1}=\bm u_k+\Delta t\, f(\bm u_k)$.''
\end{reply}
\item[p. 12] ``among these two choices'' $\rightarrow$ ``between these two choices''
\begin{reply}
Thank you for spotting the typo; we corrected it as suggested.
\end{reply}
\end{enumerate}
\end{quoteA}

\reviewerB{\section{Reviewer \#2}}

We thank the reviewer for the comments, which we address below.

\begin{quoteB}
\begin{enumerate}
\item[Sec. 3] [Comment quoted verbatim.]
\begin{reply}
[Reply, and where the manuscript changed.]
\end{reply}
\end{enumerate}
\end{quoteB}

\end{document}
